% ============================================================
\chapter*{Preface}
\addcontentsline{toc}{chapter}{Preface}
% ============================================================

\begin{center}
\textit{``All models are wrong, but some are useful.''}\\[4pt]
--- George E.\ P.\ Box
\end{center}

\bigskip

Mathematical modelling is the art and science of translating real-world
phenomena into rigorous mathematical structures, analysing those structures
to derive predictions, and confronting those predictions with observations.
This course is aimed at second- and third-year undergraduate students (L2/L3)
in mathematics, physics, quantitative biology, or engineering, with a solid
background in calculus, linear algebra, and ordinary differential equations.

\section*{Learning Objectives}

By the end of this course, the student will be able to:

\begin{enumerate}
  \item \textbf{Formulate} a concrete problem as a mathematical model,
        identifying relevant variables, parameters, and assumptions.
  \item \textbf{Analyse} the resulting equations qualitatively and
        quantitatively: equilibria, stability, bifurcations, asymptotic
        behaviour.
  \item \textbf{Simulate} models numerically using Python
        (\texttt{scipy}, \texttt{numpy}, \texttt{matplotlib}).
  \item \textbf{Validate} a model by comparing it with experimental data
        and discussing its limitations.
  \item \textbf{Communicate} results clearly and rigorously, both in
        writing and through visualisations.
\end{enumerate}

\section*{Course Philosophy}

We follow a \emph{concrete-to-abstract} approach.  Each chapter starts from
an observable phenomenon --- population growth, epidemic spread, heat
diffusion, price fluctuation --- and builds the corresponding mathematical
model step by step.  Theoretical tools (dimensional analysis, stability
theory, stochastic processes) are introduced \emph{when they become
necessary}, not as abstract prerequisites.

\begin{warning}[title=\textbf{Model $\neq$ Reality}]
A model is always a simplification.  Box's quote above should be kept in
mind throughout the course.  We will systematically emphasise the
\textbf{assumptions} underlying each model and the consequences of
violating them.
\end{warning}

\section*{Course Structure}

The course is organised into ten chapters, following a logical progression:

\begin{description}
  \item[Chapter 1 --- The Modelling Process.]
    General methodology, stages of modelling, classification of models.
  \item[Chapter 2 --- Dimensional Analysis and Scaling.]
    Buckingham $\Pi$ theorem, non-dimensionalisation, orders of magnitude.
  \item[Chapter 3 --- Population Dynamics.]
    Malthus, Verhulst, Lotka--Volterra models; stability and phase portraits.
  \item[Chapter 4 --- Epidemiological Models.]
    SIR, SEIR; basic reproduction number $\mathcal{R}_0$; vaccination
    strategies.
  \item[Chapter 5 --- Discrete Dynamical Systems and Chaos.]
    Logistic map, Lyapunov exponents, bifurcation diagrams.
  \item[Chapter 6 --- Diffusion and Transport Models.]
    Heat equation, advection-diffusion, Fick's law.
  \item[Chapter 7 --- Economic and Financial Models.]
    Solow model, Black--Scholes, supply-demand equilibrium.
  \item[Chapter 8 --- Optimization and Variational Models.]
    Linear programming, KKT conditions, calculus of variations.
  \item[Chapter 9 --- Stochastic Models.]
    Random walks, Markov chains, Brownian motion, SDEs.
  \item[Chapter 10 --- Case Studies and Projects.]
    Integrative applications spanning multiple domains.
\end{description}

\section*{Prerequisites}

\begin{itemize}
  \item \textbf{Calculus:} sequences, series, differential and integral
        calculus (single and multivariable), ordinary differential equations.
  \item \textbf{Linear Algebra:} vector spaces, matrices, eigenvalues,
        diagonalisation.
  \item \textbf{Probability:} probability spaces, random variables,
        expectation, variance, normal distribution.
  \item \textbf{Programming:} basic Python; additional libraries will be
        introduced as needed.
\end{itemize}

\section*{Conventions and Notation}

\begin{itemize}
  \item $\R$, $\N$, $\mathbb{Z}$, $\mathbb{Q}$, $\mathbb{C}$: standard
        number sets.
  \item $\E[X]$, $\Var(X)$: expectation and variance of a random
        variable~$X$.
  \item $\norm{\cdot}$, $\abs{\cdot}$: norm and absolute value.
  \item Vectors in boldface: $\mathbf{x}$, $\mathbf{v}$.
  \item Time derivatives: $\dot{x} = dx/dt$.
  \item Dimensionless quantities: $\tilde{x} = x/x_0$.
\end{itemize}

\section*{Practical Organisation}

Each chapter contains:
\begin{itemize}
  \item Numbered \textbf{definitions}, \textbf{theorems}, and
        \textbf{propositions}.
  \item Detailed \textbf{examples} illustrating the concepts.
  \item Boxed \textbf{Key Formulas}, \textbf{Warnings},
        \textbf{Best Practices}, and \textbf{Intuition} panels.
  \item Complete, commented \textbf{Python simulations}.
  \item \textbf{Exercises} of increasing difficulty.
\end{itemize}

\section*{Software}

\begin{codeblock}[title=\textbf{Python Libraries Used}]
\begin{minted}{python}
import numpy as np
import matplotlib.pyplot as plt
from scipy.integrate import odeint, solve_ivp
from scipy.optimize import minimize, linprog
from scipy.linalg import eig
\end{minted}
\end{codeblock}

\section*{Acknowledgements}

This course draws on the classic textbooks by C.~Lin and L.~Segel
(\textit{Mathematics Applied to Deterministic Problems in the Natural
Sciences}), J.\,D.~Murray (\textit{Mathematical Biology}),
S.\,H.~Strogatz (\textit{Nonlinear Dynamics and Chaos}), and
E.\,A.~Bender (\textit{An Introduction to Mathematical Modeling}).

\bigskip
\begin{flushright}
\textit{Happy reading and happy modelling!}
\end{flushright}
